%% ServeSID -- WWW submission -- Supplementary Material
%% Compile together with main.tex in the same folder.

\documentclass[sigconf, anonymous]{acmart}

\usepackage{amsmath}
\usepackage{amssymb}
\usepackage{booktabs}
\usepackage{xcolor}

\newcommand{\TODO}[1]{\textcolor{red}{\textbf{[TODO: #1]}}}

\newcommand{\method}{\textsc{ServeSID}}

\settopmatter{printacmref=false, printccs=false, printfolios=true}

\title{\method: Hierarchy-Aligned and Decision-Attributed Knowledge
Distillation for Semantic-ID Generative Recommendation\\
Supplementary Material}

\author{Anonymous Author(s)}
\affiliation{%
  \institution{Anonymous Institution}
  \city{Anonymous City}
  \country{Anonymous Country}
}
\email{anonymous@example.com}

\renewcommand{\shortauthors}{Anonymous}

\begin{document}

\maketitle

\section{Dataset Statistics}
\label{sec:sup-datasets}

\TODO{fill in the Beauty/Toys/Video Games statistics (users, items,
interactions, average sequence length, number of valid SIDs) in the
same style as the previous submission.}

\begin{table}[h]
\centering
\small
\begin{tabular}{lcccc}
\toprule
Dataset & \#Users & \#Items & \#Interactions & Avg.\ len. \\
\midrule
Amazon Beauty 2014 & \TODO{--} & \TODO{--} & \TODO{--} & \TODO{--} \\
Amazon Toys 2014   & \TODO{--} & \TODO{--} & \TODO{--} & \TODO{--} \\
Video Games 2023   & \TODO{--} & \TODO{--} & \TODO{--} & \TODO{--} \\
\bottomrule
\end{tabular}
\caption{Dataset statistics.}
\label{tab:sup-stats}
\end{table}

\section{Implementation and Hyperparameters}
\label{sec:sup-hyper}

All methods share the same fixed TIGER teacher checkpoint, RQ-VAE
tokenizer, and four-token SID space with codebook size 256. The
three student backbones encode the last 20 interactions and score
each SID depth with lightweight heads:

\begin{itemize}
\item \textbf{MLP}: 1.664M trainable parameters;
\item \textbf{GRU}: hidden size 384, 1.201M trainable parameters;
\item \textbf{CNN}: embedding dimension 256, 1.753M trainable parameters.
\end{itemize}

\noindent\emph{Caveat:} the GRU/CNN runs used slightly different
parameter budgets than the MLP due to a launch-script issue; a
strictly matched rerun is planned and will be reported in the final
version.

Students decode with constrained beam search of width $W=50$
(top-50); the TIGER teacher reference is decoded with beam 20 (the
genrec default) and is reported for scale only, since the absolute
values are not directly comparable across beam widths.

ServeSID hyperparameters: support size $M=64$, training cutoff
$K_{\mathrm{tr}}=10$, maximum outsider rank $B=40$, online beam
$W=50$, HServe weight $\lambda_H=4$, DA-CDRS weight $\lambda_D=1$,
depth weights $\alpha_d=1$ for $d=1,\dots,L$, pair margin $\mu=.05$,
increment-correction weight $\lambda_{\mathrm{corr}}=31.42$, learning
rate $5\times10^{-5}$, seeds $\{42,43,44\}$. The draft tables have not yet been verified
against the no-direct-pair-anchor checkpoints and should not be
interpreted as validated three-seed means. Teacher supports, marginals, pairs, and increment
scales are built from training requests only; validation/test teacher
outputs never enter optimization.

\section{Full Main-Result Tables}
\label{sec:sup-main}

Tables~\ref{tab:supp-utility-Beauty}--\ref{tab:supp-jac-VG} give the
full per-dataset metrics (R@5, N@5, R@10, N@10) and teacher-agreement
(Jac@5, Jac@10) for all methods and students.

\begin{table*}[t]
\centering
\footnotesize
\setlength{\tabcolsep}{1.1mm}
\renewcommand{\arraystretch}{1.05}
\begin{tabular}{lcccccccccccc}
\toprule
Method & \multicolumn{4}{c}{MLP} & \multicolumn{4}{c}{GRU} & \multicolumn{4}{c}{CNN} \\
 & R@5 & N@5 & R@10 & N@10 & R@5 & N@5 & R@10 & N@10 & R@5 & N@5 & R@10 & N@10 \\
\cmidrule(lr){2-5}\cmidrule(lr){6-9}\cmidrule(lr){10-13}
\midrule
\multicolumn{13}{l}{\emph{Teacher TIGER reference (beam-20):}} \\
TIGER & 0.043 & 0.028 & 0.068 & 0.036 & 0.043 & 0.028 & 0.068 & 0.036 & 0.043 & 0.028 & 0.068 & 0.036 \\
\midrule
SID-MLP (base KD) & 0.0429 & 0.0282 & 0.0660 & 0.0356 & 0.0416 & 0.0281 & 0.0656 & 0.0358 & 0.0380 & 0.0254 & 0.0580 & 0.0319 \\
SmartGR & 0.0434 & 0.0286 & 0.0668 & 0.0363 & 0.0426 & 0.0286 & 0.0669 & 0.0361 & 0.0374 & 0.0250 & 0.0594 & 0.0320 \\
RD & 0.0422 & 0.0281 & 0.0657 & 0.0357 & 0.0433 & 0.0285 & 0.0671 & 0.0362 & 0.0419 & 0.0274 & 0.0649 & 0.0348 \\
RRD & 0.0430 & 0.0287 & 0.0659 & 0.0361 & 0.0422 & 0.0281 & 0.0664 & 0.0359 & 0.0420 & 0.0281 & 0.0649 & 0.0354 \\
RCE-KD & 0.0434 & 0.0286 & 0.0666 & 0.0363 & 0.0432 & 0.0286 & 0.0671 & 0.0361 & 0.0409 & 0.0276 & 0.0651 & 0.0354 \\
GKD & 0.0428 & 0.0286 & 0.0668 & 0.0363 & 0.0416 & 0.0281 & 0.0656 & 0.0358 & 0.0380 & 0.0254 & 0.0580 & 0.0319 \\
DistiLLM-2 & 0.0428 & 0.0287 & 0.0666 & 0.0361 & 0.0416 & 0.0281 & 0.0656 & 0.0358 & 0.0380 & 0.0254 & 0.0580 & 0.0319 \\
\textbf{\method} & 0.0442 & 0.0292 & 0.0679 & 0.0369 & 0.0441 & 0.0291 & 0.0682 & 0.0368 & 0.0437 & 0.0288 & 0.0686 & 0.0368 \\
\bottomrule
\end{tabular}
\caption{Full utility metrics on Beauty for all methods and students.}
\label{tab:supp-utility-Beauty}
\end{table*}

\begin{table*}[t]
\centering
\footnotesize
\setlength{\tabcolsep}{2mm}
\renewcommand{\arraystretch}{1.05}
\begin{tabular}{lcccccc}
\toprule
Method & \multicolumn{2}{c}{MLP} & \multicolumn{2}{c}{GRU} & \multicolumn{2}{c}{CNN} \\
 & Jac@5 & Jac@10 & Jac@5 & Jac@10 & Jac@5 & Jac@10 \\
\cmidrule(lr){2-3}\cmidrule(lr){4-5}\cmidrule(lr){6-7}
\midrule
SID-MLP (base KD) & 0.0113 & 0.0154 & 0.0091 & 0.0122 & 0.0075 & 0.0103 \\
SmartGR & 0.0115 & 0.0155 & 0.0092 & 0.0125 & 0.0075 & 0.0106 \\
RD & 0.0098 & 0.0141 & 0.0096 & 0.0138 & 0.0092 & 0.0135 \\
RRD & 0.0125 & 0.0154 & 0.0116 & 0.0143 & 0.0118 & 0.0145 \\
RCE-KD & 0.0112 & 0.0151 & 0.0104 & 0.0140 & 0.0094 & 0.0124 \\
GKD & 0.0115 & 0.0157 & 0.0091 & 0.0122 & 0.0075 & 0.0103 \\
DistiLLM-2 & 0.0117 & 0.0158 & 0.0091 & 0.0122 & 0.0075 & 0.0103 \\
\textbf{\method} & 0.0118 & 0.0154 & 0.0116 & 0.0151 & 0.0114 & 0.0150 \\
\bottomrule
\end{tabular}
\caption{Teacher-agreement (Jaccard with the TIGER beam-20 top-K list) on Beauty for all methods and students.}
\label{tab:supp-jac-Beauty}
\end{table*}

\begin{table*}[t]
\centering
\footnotesize
\setlength{\tabcolsep}{1.1mm}
\renewcommand{\arraystretch}{1.05}
\begin{tabular}{lcccccccccccc}
\toprule
Method & \multicolumn{4}{c}{MLP} & \multicolumn{4}{c}{GRU} & \multicolumn{4}{c}{CNN} \\
 & R@5 & N@5 & R@10 & N@10 & R@5 & N@5 & R@10 & N@10 & R@5 & N@5 & R@10 & N@10 \\
\cmidrule(lr){2-5}\cmidrule(lr){6-9}\cmidrule(lr){10-13}
\midrule
\multicolumn{13}{l}{\emph{Teacher TIGER reference (beam-20):}} \\
TIGER & 0.041 & 0.025 & 0.065 & 0.033 & 0.041 & 0.025 & 0.065 & 0.033 & 0.041 & 0.025 & 0.065 & 0.033 \\
\midrule
SID-MLP (base KD) & 0.0374 & 0.0235 & 0.0626 & 0.0315 & 0.0367 & 0.0233 & 0.0565 & 0.0296 & 0.0328 & 0.0210 & 0.0533 & 0.0276 \\
SmartGR & 0.0381 & 0.0241 & 0.0623 & 0.0319 & 0.0382 & 0.0242 & 0.0604 & 0.0314 & 0.0329 & 0.0206 & 0.0520 & 0.0267 \\
RD & 0.0380 & 0.0234 & 0.0615 & 0.0309 & 0.0381 & 0.0245 & 0.0614 & 0.0318 & 0.0380 & 0.0241 & 0.0603 & 0.0312 \\
RRD & 0.0396 & 0.0245 & 0.0625 & 0.0319 & 0.0381 & 0.0245 & 0.0622 & 0.0322 & 0.0378 & 0.0239 & 0.0585 & 0.0306 \\
RCE-KD & 0.0387 & 0.0244 & 0.0631 & 0.0324 & 0.0380 & 0.0246 & 0.0622 & 0.0324 & 0.0364 & 0.0232 & 0.0587 & 0.0303 \\
GKD & 0.0373 & 0.0237 & 0.0627 & 0.0318 & 0.0382 & 0.0243 & 0.0613 & 0.0316 & 0.0328 & 0.0210 & 0.0533 & 0.0276 \\
DistiLLM-2 & 0.0375 & 0.0238 & 0.0614 & 0.0315 & 0.0367 & 0.0233 & 0.0565 & 0.0296 & 0.0328 & 0.0210 & 0.0533 & 0.0276 \\
\textbf{\method} & 0.0408 & 0.0254 & 0.0644 & 0.0330 & 0.0388 & 0.0250 & 0.0650 & 0.0334 & 0.0398 & 0.0250 & 0.0639 & 0.0328 \\
\bottomrule
\end{tabular}
\caption{Full utility metrics on Toys for all methods and students.}
\label{tab:supp-utility-Toys}
\end{table*}

\begin{table*}[t]
\centering
\footnotesize
\setlength{\tabcolsep}{2mm}
\renewcommand{\arraystretch}{1.05}
\begin{tabular}{lcccccc}
\toprule
Method & \multicolumn{2}{c}{MLP} & \multicolumn{2}{c}{GRU} & \multicolumn{2}{c}{CNN} \\
 & Jac@5 & Jac@10 & Jac@5 & Jac@10 & Jac@5 & Jac@10 \\
\cmidrule(lr){2-3}\cmidrule(lr){4-5}\cmidrule(lr){6-7}
\midrule
SID-MLP (base KD) & 0.0062 & 0.0089 & 0.0045 & 0.0066 & 0.0038 & 0.0056 \\
SmartGR & 0.0065 & 0.0089 & 0.0045 & 0.0068 & 0.0041 & 0.0061 \\
RD & 0.0054 & 0.0080 & 0.0051 & 0.0077 & 0.0047 & 0.0076 \\
RRD & 0.0061 & 0.0084 & 0.0057 & 0.0076 & 0.0056 & 0.0080 \\
RCE-KD & 0.0061 & 0.0087 & 0.0052 & 0.0075 & 0.0043 & 0.0066 \\
GKD & 0.0065 & 0.0091 & 0.0048 & 0.0072 & 0.0038 & 0.0056 \\
DistiLLM-2 & 0.0066 & 0.0092 & 0.0045 & 0.0066 & 0.0038 & 0.0056 \\
\textbf{\method} & 0.0062 & 0.0087 & 0.0060 & 0.0086 & 0.0059 & 0.0083 \\
\bottomrule
\end{tabular}
\caption{Teacher-agreement (Jaccard with the TIGER beam-20 top-K list) on Toys for all methods and students.}
\label{tab:supp-jac-Toys}
\end{table*}

\begin{table*}[t]
\centering
\footnotesize
\setlength{\tabcolsep}{1.1mm}
\renewcommand{\arraystretch}{1.05}
\begin{tabular}{lcccccccccccc}
\toprule
Method & \multicolumn{4}{c}{MLP} & \multicolumn{4}{c}{GRU} & \multicolumn{4}{c}{CNN} \\
 & R@5 & N@5 & R@10 & N@10 & R@5 & N@5 & R@10 & N@10 & R@5 & N@5 & R@10 & N@10 \\
\cmidrule(lr){2-5}\cmidrule(lr){6-9}\cmidrule(lr){10-13}
\midrule
\multicolumn{13}{l}{\emph{Teacher TIGER reference (beam-20):}} \\
TIGER & 0.044 & 0.028 & 0.070 & 0.036 & 0.044 & 0.028 & 0.070 & 0.036 & 0.044 & 0.028 & 0.070 & 0.036 \\
\midrule
SID-MLP (base KD) & 0.0588 & 0.0383 & 0.0922 & 0.0488 & 0.0591 & 0.0389 & 0.0935 & 0.0499 & 0.0553 & 0.0358 & 0.0870 & 0.0460 \\
SmartGR & 0.0593 & 0.0388 & 0.0929 & 0.0496 & 0.0594 & 0.0391 & 0.0933 & 0.0500 & 0.0549 & 0.0359 & 0.0870 & 0.0462 \\
RD & 0.0575 & 0.0372 & 0.0927 & 0.0485 & 0.0595 & 0.0388 & 0.0938 & 0.0498 & 0.0580 & 0.0373 & 0.0921 & 0.0483 \\
RRD & 0.0592 & 0.0390 & 0.0918 & 0.0496 & 0.0596 & 0.0393 & 0.0892 & 0.0490 & 0.0591 & 0.0389 & 0.0884 & 0.0483 \\
RCE-KD & 0.0593 & 0.0388 & 0.0926 & 0.0497 & 0.0595 & 0.0394 & 0.0938 & 0.0503 & 0.0589 & 0.0387 & 0.0919 & 0.0493 \\
GKD & 0.0593 & 0.0388 & 0.0929 & 0.0496 & 0.0597 & 0.0393 & 0.0937 & 0.0502 & 0.0569 & 0.0370 & 0.0896 & 0.0475 \\
DistiLLM-2 & 0.0594 & 0.0387 & 0.0927 & 0.0494 & 0.0596 & 0.0392 & 0.0939 & 0.0503 & 0.0587 & 0.0384 & 0.0914 & 0.0489 \\
\textbf{\method} & 0.0604 & 0.0396 & 0.0944 & 0.0505 & 0.0607 & 0.0400 & 0.0954 & 0.0512 & 0.0605 & 0.0399 & 0.0944 & 0.0508 \\
\bottomrule
\end{tabular}
\caption{Full utility metrics on VG for all methods and students.}
\label{tab:supp-utility-VG}
\end{table*}

\begin{table*}[t]
\centering
\footnotesize
\setlength{\tabcolsep}{2mm}
\renewcommand{\arraystretch}{1.05}
\begin{tabular}{lcccccc}
\toprule
Method & \multicolumn{2}{c}{MLP} & \multicolumn{2}{c}{GRU} & \multicolumn{2}{c}{CNN} \\
 & Jac@5 & Jac@10 & Jac@5 & Jac@10 & Jac@5 & Jac@10 \\
\cmidrule(lr){2-3}\cmidrule(lr){4-5}\cmidrule(lr){6-7}
\midrule
SID-MLP (base KD) & 0.0100 & 0.0131 & 0.0088 & 0.0117 & 0.0092 & 0.0121 \\
SmartGR & 0.0100 & 0.0130 & 0.0089 & 0.0118 & 0.0091 & 0.0121 \\
RD & 0.0086 & 0.0116 & 0.0083 & 0.0113 & 0.0081 & 0.0112 \\
RRD & 0.0094 & 0.0114 & 0.0090 & 0.0105 & 0.0090 & 0.0109 \\
RCE-KD & 0.0097 & 0.0126 & 0.0091 & 0.0119 & 0.0092 & 0.0120 \\
GKD & 0.0101 & 0.0131 & 0.0089 & 0.0118 & 0.0093 & 0.0124 \\
DistiLLM-2 & 0.0100 & 0.0131 & 0.0093 & 0.0121 & 0.0097 & 0.0127 \\
\textbf{\method} & 0.0094 & 0.0122 & 0.0092 & 0.0121 & 0.0092 & 0.0121 \\
\bottomrule
\end{tabular}
\caption{Teacher-agreement (Jaccard with the TIGER beam-20 top-K list) on VG for all methods and students.}
\label{tab:supp-jac-VG}
\end{table*}


\section{Full Ablation Tables}
\label{sec:sup-ablation}

Tables~\ref{tab:supp-ablation-Beauty}--\ref{tab:supp-ablation-VG} give
all six metrics for every ablation arm on each dataset. Ablation arms
are defined in the main paper (Section~6.3).

\begin{table*}[t]
\centering
\footnotesize
\setlength{\tabcolsep}{1.1mm}
\renewcommand{\arraystretch}{1.05}
\begin{tabular}{lcccccccccccc}
\toprule
Arm & \multicolumn{4}{c}{MLP} & \multicolumn{4}{c}{GRU} & \multicolumn{4}{c}{CNN} \\
 & R@5 & N@5 & R@10 & N@10 & R@5 & N@5 & R@10 & N@10 & R@5 & N@5 & R@10 & N@10 \\
\cmidrule(lr){2-5}\cmidrule(lr){6-9}\cmidrule(lr){10-13}
\midrule
SID-MLP (base KD) & 0.0429 & 0.0282 & 0.0660 & 0.0356 & 0.0416 & 0.0281 & 0.0656 & 0.0358 & 0.0380 & 0.0254 & 0.0580 & 0.0319 \\
d4-only KL & 0.0431 & 0.0283 & 0.0663 & 0.0359 & 0.0422 & 0.0283 & 0.0660 & 0.0360 & 0.0402 & 0.0266 & 0.0625 & 0.0348 \\
DA-CDRS only & 0.0435 & 0.0286 & 0.0672 & 0.0362 & 0.0430 & 0.0286 & 0.0673 & 0.0364 & 0.0414 & 0.0273 & 0.0652 & 0.0359 \\
HServe only & 0.0436 & 0.0287 & 0.0670 & 0.0363 & 0.0431 & 0.0286 & 0.0671 & 0.0364 & 0.0418 & 0.0276 & 0.0656 & 0.0360 \\
HServe + Uniform & 0.0437 & 0.0286 & 0.0673 & 0.0364 & 0.0437 & 0.0287 & 0.0674 & 0.0365 & 0.0423 & 0.0280 & 0.0669 & 0.0362 \\
\textbf{\method (full)} & \textbf{0.0442} & \textbf{0.0292} & \textbf{0.0679} & \textbf{0.0369} & \textbf{0.0441} & \textbf{0.0291} & \textbf{0.0682} & \textbf{0.0368} & \textbf{0.0437} & \textbf{0.0288} & \textbf{0.0686} & \textbf{0.0368} \\
\bottomrule
\end{tabular}
\caption{Full ablation metrics on Beauty for all students.}
\label{tab:supp-ablation-Beauty}
\end{table*}

\begin{table*}[t]
\centering
\footnotesize
\setlength{\tabcolsep}{1.1mm}
\renewcommand{\arraystretch}{1.05}
\begin{tabular}{lcccccccccccc}
\toprule
Arm & \multicolumn{4}{c}{MLP} & \multicolumn{4}{c}{GRU} & \multicolumn{4}{c}{CNN} \\
 & R@5 & N@5 & R@10 & N@10 & R@5 & N@5 & R@10 & N@10 & R@5 & N@5 & R@10 & N@10 \\
\cmidrule(lr){2-5}\cmidrule(lr){6-9}\cmidrule(lr){10-13}
\midrule
SID-MLP (base KD) & 0.0374 & 0.0235 & 0.0626 & 0.0315 & 0.0367 & 0.0233 & 0.0565 & 0.0296 & 0.0328 & 0.0210 & 0.0533 & 0.0276 \\
d4-only KL & 0.0379 & 0.0237 & 0.0628 & 0.0317 & 0.0369 & 0.0237 & 0.0594 & 0.0313 & 0.0354 & 0.0225 & 0.0566 & 0.0305 \\
DA-CDRS only & 0.0389 & 0.0246 & 0.0638 & 0.0324 & 0.0377 & 0.0241 & 0.0622 & 0.0323 & 0.0370 & 0.0236 & 0.0598 & 0.0322 \\
HServe only & 0.0390 & 0.0245 & 0.0633 & 0.0323 & 0.0379 & 0.0243 & 0.0630 & 0.0325 & 0.0375 & 0.0238 & 0.0606 & 0.0320 \\
HServe + Uniform & 0.0395 & 0.0248 & 0.0640 & 0.0325 & 0.0381 & 0.0247 & 0.0642 & 0.0328 & 0.0382 & 0.0242 & 0.0630 & 0.0323 \\
\textbf{\method (full)} & \textbf{0.0408} & \textbf{0.0254} & \textbf{0.0644} & \textbf{0.0330} & \textbf{0.0388} & \textbf{0.0250} & \textbf{0.0650} & \textbf{0.0334} & \textbf{0.0398} & \textbf{0.0250} & \textbf{0.0639} & \textbf{0.0328} \\
\bottomrule
\end{tabular}
\caption{Full ablation metrics on Toys for all students.}
\label{tab:supp-ablation-Toys}
\end{table*}

\begin{table*}[t]
\centering
\footnotesize
\setlength{\tabcolsep}{1.1mm}
\renewcommand{\arraystretch}{1.05}
\begin{tabular}{lcccccccccccc}
\toprule
Arm & \multicolumn{4}{c}{MLP} & \multicolumn{4}{c}{GRU} & \multicolumn{4}{c}{CNN} \\
 & R@5 & N@5 & R@10 & N@10 & R@5 & N@5 & R@10 & N@10 & R@5 & N@5 & R@10 & N@10 \\
\cmidrule(lr){2-5}\cmidrule(lr){6-9}\cmidrule(lr){10-13}
\midrule
SID-MLP (base KD) & 0.0588 & 0.0383 & 0.0922 & 0.0488 & 0.0591 & 0.0389 & 0.0935 & 0.0499 & 0.0553 & 0.0358 & 0.0870 & 0.0460 \\
d4-only KL & 0.0589 & 0.0384 & 0.0925 & 0.0491 & 0.0593 & 0.0390 & 0.0937 & 0.0502 & 0.0570 & 0.0370 & 0.0888 & 0.0488 \\
DA-CDRS only & 0.0595 & 0.0388 & 0.0932 & 0.0498 & 0.0601 & 0.0394 & 0.0943 & 0.0507 & 0.0586 & 0.0380 & 0.0914 & 0.0498 \\
HServe only & 0.0597 & 0.0390 & 0.0935 & 0.0499 & 0.0600 & 0.0395 & 0.0945 & 0.0506 & 0.0591 & 0.0386 & 0.0923 & 0.0500 \\
HServe + Uniform & 0.0599 & 0.0391 & 0.0938 & 0.0500 & 0.0602 & 0.0397 & 0.0950 & 0.0510 & 0.0596 & 0.0390 & 0.0935 & 0.0502 \\
\textbf{\method (full)} & \textbf{0.0604} & \textbf{0.0396} & \textbf{0.0944} & \textbf{0.0505} & \textbf{0.0607} & \textbf{0.0400} & \textbf{0.0954} & \textbf{0.0512} & \textbf{0.0605} & \textbf{0.0399} & \textbf{0.0944} & \textbf{0.0508} \\
\bottomrule
\end{tabular}
\caption{Full ablation metrics on VG for all students.}
\label{tab:supp-ablation-VG}
\end{table*}


\section{Diagnostic Metrics and Protocol}
\label{sec:sup-diagnostics}

We evaluate diagnostics on 10,000 held-out requests per domain
using teacher-proposed top-40 valid SIDs as a common support.
For each prefix depth, preferred-prefix disagreement is the fraction
of requests where teacher and student select different groups of
maximum aggregated probability. Mixed-sign pairs have at least one
positive and one negative residual among eligible depths from the
first SID divergence through $L-1$. We also report the fractions with
at least two positive residuals and with a negative largest-magnitude
residual. The cancellation ratio summarizes cancellation between
positive and negative signed depth residuals; the exact aggregation,
zero-residual handling, and tie-breaking protocol must be checked
against the diagnostic implementation before final submission.

\section{Reproducibility Notes}
\label{sec:sup-repro}

\begin{itemize}
\item Evaluation: next-item utility is Recall@$K$ and NDCG@$K$ under
leave-one-out on the held-out interaction; teacher-agreement Jac@$K$
is the Jaccard between the student's beam-50 top-$K$ list and the
teacher's beam-20 top-$K$ list.
\item The base student (SID-MLP / base KD) minimizes ground-truth SID
cross-entropy plus gold-prefix token KL (Eq.~2 of the main paper).
ServeSID replaces token KL with HServe and DA-CDRS. The
complete-pair Softplus function is used only to calculate detached
DA attribution weights; no direct pairwise anchor loss is optimized.
\item Held-out diagnostic pairs (teacher ranks $1{:}K$ vs.\
$K{+}1{:}B$) are used only for mechanism statistics and never enter
optimization.
\end{itemize}

\end{document}
